\begin{mainquestion}{mosfetprocess}{MOSFET Process Technology}

\questionpoints Consider an NMOS transistor fabricated using silicon dioxide (SiO$_2$) gate oxide 
with relative permittivity $\varepsilon_r = 3.9$.

Process parameters:
\[
L_{\min} = 0.15\,\mu m, \quad
t_{ox} = 2.7\,nm, \quad
\mu_n = 400\,\mathrm{cm^2/V\cdot s}, \quad
V_t = 0.5\,V
\]

The oxide capacitance per unit area is defined as
\[
C_{ox} = \frac{\varepsilon_{ox}}{t_{ox}}, \quad
\varepsilon_{ox} = \varepsilon_r \varepsilon_0
\]

The long–channel saturation current equation is
\[
I_D=\frac{1}{2}k_n'\frac{W}{L}(V_{GS}-V_T)^2
\]
where the process transconductance parameter is
\[
k_n' \equiv \mu_n C_{ox}
\]

%-----------------------------------------

% =========================================================
\begin{partquestion}{findcox}
{Find the oxide capacitance per unit area $C_{ox}$.}

\begin{answer}
\[
\varepsilon_{ox} = 3.9(8.854\times10^{-12}) 
= 3.45\times10^{-11}\,F/m
\]

\[
C_{ox} = \frac{\varepsilon_{ox}}{t_{ox}}
= \frac{3.45\times10^{-11}}{2.7\times10^{-9}}
= 1.28\times10^{-2}\,F/m^2
\]

\[
C_{ox} = 12.8\,fF/\mu m^2
\]
\end{answer}

%-----------------------------------------
\end{partquestion}

% =========================================================
\begin{partquestion}{findknprime}
{Find the process transconductance parameter $k_n'$.}

\begin{answer}
\[
k_n'=\mu_n C_{ox}
\]

Convert mobility:

\[
400\,cm^2/V\cdot s = 400\times10^{-4}\,m^2/V\cdot s
\]

\[
k_n'=(400\times10^{-4})(1.28\times10^{-2})
=5.12\times10^{-4}\,A/V^2
\]

\[
k_n' = 512\,\mu A/V^2
\]
\end{answer}

%-----------------------------------------
\pagebreak
\end{partquestion}

% =========================================================
\begin{partquestion}{findwidth}
{Find the channel width $W$ if the channel resistance is $1.7\,k\Omega$ at $V_{GS}=0.8\,V$ and small $V_{DS}$.

For small $V_{DS}$ (linear region), the small-signal drain conductance is
\[
g_{DS}=k_n'\frac{W}{L}(V_{GS}-V_t)
\]
and
\[
r_{DS}=\frac{1}{g_{DS}}
\]}

\begin{answer}
\[
W=\frac{L}{r_{DS}k_n'(V_{GS}-V_t)}
\]

\[
W=\frac{0.15\times10^{-6}}
{(1700)(5.12\times10^{-4})(0.8-0.5)}
\]

\[
W=5.74\times10^{-7}\,m
\]

\[
W=0.574\,\mu m
\]
\end{answer}

%-----------------------------------------
\end{partquestion}

% =========================================================
\begin{partquestion}{saturation}
{The transistor must operate in saturation with drain current $I_D=60\,\mu A$.  
Find the required gate voltage $V_{GS}$.}

\begin{answer}
\[
I_D=\frac{1}{2}k_n'\frac{W}{L}(V_{GS}-V_T)^2
\]

\[
V_{OV}=V_{GS}-V_T
=\sqrt{\frac{2I_DL}{k_n'W}}
\]

\[
V_{OV}=\sqrt{
\frac{2(60\times10^{-6})(0.15)}
{(5.12\times10^{-4})(0.574)}
}
=0.25\,V
\]

\[
V_{GS}=V_T+V_{OV}
=0.5+0.25
=0.75\,V
\]
\end{answer}

%-----------------------------------------
\end{partquestion}

% =========================================================
\begin{partquestion}{findvdsmin}
{What minimum $V_{DS}$ ensures the device remains in saturation?}

\begin{answer}
Saturation condition:

\[
V_{DS(min)} = V_{OV}
\]

\[
V_{DS(min)}=0.25\,V
\]
\end{answer}

\end{partquestion}

\end{mainquestion}
